% One pre-norm ViT block with the perturbation sites. Every box has the same
% width so the residual skip, drawn left of the widest box, crosses none of them.
% A and D sit on the branch after the skip splits off, B on the branch output
% before the addition and C on the stream after each addition, which is where
% PSBD-RD applies dropout.
\begin{tikzpicture}[
  box/.style={draw, rounded corners=1pt, minimum width=30mm, text width=28mm, minimum height=5.5mm, font=\scriptsize, align=center},
  norm/.style={box, fill=black!6},
  op/.style={box, fill=blue!7},
  plus/.style={draw, circle, inner sep=0.8pt, font=\scriptsize, fill=white},
  site/.style={circle, fill=black!70, text=white, inner sep=0pt, minimum size=3.4mm, font=\bfseries\tiny},
  leader/.style={black!45, thin},
  flow/.style={-{Latex[length=1.4mm]}, thick}
]
\def\gap{5mm}
\def\skipx{-20mm}
\def\sitex{21mm}
\node[font=\scriptsize] (in) at (0,0) {residual stream, \VitTokens{} tokens $\times$ \VitWidth{}};
\coordinate[below=\gap of in] (split1);
\node[norm, below=\gap of split1] (ln1) {LayerNorm};
\node[op, below=\gap of ln1] (attn) {Multi-Head Attention};
\node[plus, below=\gap of attn] (add1) {$+$};
\coordinate[below=\gap of add1] (split2);
\node[norm, below=\gap of split2] (ln2) {LayerNorm};
\node[op, below=\gap of ln2] (mlp) {MLP, per token};
\node[plus, below=\gap of mlp] (add2) {$+$};
\node[font=\scriptsize, below=\gap of add2] (out) {to the next block};
\draw[flow] (in) -- (ln1);
\draw[flow] (ln1) -- (attn);
\draw[flow] (attn) -- (add1);
\draw[flow] (add1) -- (ln2);
\draw[flow] (ln2) -- (mlp);
\draw[flow] (mlp) -- (add2);
\draw[flow] (add2) -- (out);
\draw[flow] (split1) -- (\skipx,0 |- split1) |- (add1.west);
\draw[flow] (split2) -- (\skipx,0 |- split2) |- (add2.west);
\fill (split1) circle (0.5pt);
\fill (split2) circle (0.5pt);
% Each marker sits in 1 column right of the boxes, joined to the midpoint of
% the arrow it perturbs.
\coordinate (a) at ($(split1)!0.5!(ln1.north)$);
\coordinate (b) at ($(attn.south)!0.5!(add1.north)$);
\coordinate (c1) at ($(add1.south)!0.35!(split2)$);
\coordinate (d) at ($(split2)!0.5!(ln2.north)$);
\coordinate (c2) at ($(add2.south)!0.5!(out.north)$);
\foreach \point/\name in {a/A, b/B, c1/C, d/D, c2/C} {
  \draw[leader] (\point) -- (\sitex,0 |- \point);
  \node[site] at (\sitex,0 |- \point) {\name};
}
\end{tikzpicture}
